\documentclass[landscape,11pt]{article}
\usepackage[utf8]{inputenc}
\usepackage[margin=1cm]{geometry}
\usepackage{multicol}
\usepackage[dvipsnames]{xcolor}
\usepackage{helvet}
\renewcommand{\familydefault}{\sfdefault}
\usepackage{tabularx}
\usepackage{tcolorbox}
\usepackage{microtype}
\usepackage{fancyhdr}
\usepackage{hyperref}
\hypersetup{hidelinks}
\urlstyle{same}

\definecolor{footgray}{gray}{0.35}
\definecolor{claudeorange}{RGB}{217,119,87}
\definecolor{claudeblue}{RGB}{59,130,246}
\definecolor{claudegreen}{RGB}{34,197,94}
\definecolor{claudepurple}{RGB}{147,51,234}
\definecolor{claudegray}{RGB}{75,85,99}

\setlength{\parindent}{0pt}
\setlength{\columnsep}{1cm}

\pagestyle{fancy}
\fancyhf{}
\renewcommand{\headrulewidth}{0pt}
\renewcommand{\footrulewidth}{0.2pt}
\setlength{\footskip}{16pt}
\fancyfoot[R]{\footnotesize\color{footgray} \url{https://github.com/atarantino/cheatsheets}}

\begin{document}
\begin{center}
{\Huge \textbf{CLAUDE CODE CHEAT SHEET}}\\[3mm]
{\large Anthropic's AI-powered CLI for software engineering}\\[5mm]
\end{center}

\begin{multicols}{3}

\section*{\textcolor{claudeorange}{1. GETTING STARTED}}
\renewcommand{\arraystretch}{1.15}
\begin{tabularx}{\linewidth}{@{}>{\ttfamily}l>{\raggedright\arraybackslash}X@{}}
\normalfont\textbf{Command} & \textbf{Description}\\
\hline
claude & Start interactive REPL\\
claude "prompt" & Run with initial prompt\\
claude -p "prompt" & Print mode (non-interactive)\\
claude -c & Continue last conversation\\
claude --resume & Resume specific session\\
cat file | claude & Pipe input to Claude\\
\end{tabularx}
\renewcommand{\arraystretch}{1}

\vspace{2mm}
\begin{tcolorbox}[colback=claudeorange!8, colframe=claudeorange, boxsep=2pt, arc=1pt, fontupper=\small]
\textbf{Installation:} \texttt{npm install -g @anthropic-ai/claude-code}
\end{tcolorbox}

\section*{\textcolor{claudeblue}{2. SLASH COMMANDS}}
\renewcommand{\arraystretch}{1.15}
\begin{tabularx}{\linewidth}{@{}>{\ttfamily}l>{\raggedright\arraybackslash}X@{}}
\normalfont\textbf{Command} & \textbf{Action}\\
\hline
/help & Show help and commands\\
/clear & Clear conversation history\\
/compact & Summarize conversation\\
/config & Open settings\\
/cost & Show token usage\\
/doctor & Check installation health\\
/init & Initialize project config\\
/login & Authenticate with Anthropic\\
/logout & Sign out\\
/memory & Edit memory files\\
/model & Switch AI model\\
/permissions & Manage tool permissions\\
/review & Code review mode\\
/terminal-setup & Configure terminal\\
/vim & Toggle vim keybindings\\
\end{tabularx}
\renewcommand{\arraystretch}{1}

\vspace{2mm}
\begin{tcolorbox}[colback=claudeblue!8, colframe=claudeblue, boxsep=2pt, arc=1pt, fontupper=\small\itshape]
Type \texttt{/} in the REPL to see all available commands.
\end{tcolorbox}

\section*{\textcolor{claudegreen}{3. KEYBOARD SHORTCUTS}}
\renewcommand{\arraystretch}{1.15}
\begin{tabularx}{\linewidth}{@{}>{\ttfamily}l>{\raggedright\arraybackslash}X@{}}
\normalfont\textbf{Shortcut} & \textbf{Action}\\
\hline
Enter & Send message (single line)\\
Shift+Enter & New line in input\\
Ctrl+C & Cancel current operation\\
Ctrl+D & Exit Claude Code\\
Esc & Interrupt response\\
Esc, Esc & Force quit\\
Up/Down & Navigate history\\
Tab & Autocomplete files/commands\\
\end{tabularx}
\renewcommand{\arraystretch}{1}

\vspace{2mm}
\begin{tcolorbox}[colback=claudegreen!8, colframe=claudegreen, boxsep=2pt, arc=1pt, fontupper=\small\itshape]
Enable vim mode with \texttt{/vim} for modal editing.
\end{tcolorbox}

\section*{\textcolor{claudepurple}{4. CORE TOOLS}}
Claude Code uses these tools to interact with your system:

\renewcommand{\arraystretch}{1.15}
\begin{tabularx}{\linewidth}{@{}>{\ttfamily}l>{\raggedright\arraybackslash}X@{}}
\normalfont\textbf{Tool} & \textbf{Purpose}\\
\hline
Read & Read file contents\\
Edit & Make targeted edits\\
Write & Create/overwrite files\\
Bash & Execute shell commands\\
Glob & Find files by pattern\\
Grep & Search file contents\\
WebFetch & Fetch URL content\\
WebSearch & Search the web\\
Task & Launch sub-agents\\
\end{tabularx}
\renewcommand{\arraystretch}{1}

\vspace{2mm}
\begin{tcolorbox}[colback=claudepurple!8, colframe=claudepurple, boxsep=2pt, arc=1pt, fontupper=\small]
\textbf{Permission modes:}\\
\texttt{--allowedTools} — Whitelist specific tools\\
\texttt{--dangerously-skip-permissions} — Auto-approve all
\end{tcolorbox}

\section*{\textcolor{BrickRed}{5. GIT OPERATIONS}}
\renewcommand{\arraystretch}{1.15}
\begin{tabularx}{\linewidth}{@{}>{\ttfamily}l>{\raggedright\arraybackslash}X@{}}
\normalfont\textbf{Request} & \textbf{Claude Will...}\\
\hline
"commit this" & Stage, write message, commit\\
"create a PR" & Push branch, create PR via gh\\
"review this PR" & Analyze changes, give feedback\\
"git status" & Show repo state\\
"undo last commit" & Reset safely\\
\end{tabularx}
\renewcommand{\arraystretch}{1}

\vspace{2mm}
\begin{tcolorbox}[colback=BrickRed!8, colframe=BrickRed, boxsep=2pt, arc=1pt, fontupper=\small\itshape]
Claude follows git best practices and won't force-push to main.
\end{tcolorbox}

\section*{\textcolor{claudegray}{6. CONFIGURATION}}

\textbf{Config file locations:}
\begin{itemize}
  \setlength{\itemsep}{0pt}
  \item \texttt{\textasciitilde/.claude/settings.json} — Global settings
  \item \texttt{.claude/settings.json} — Project settings
  \item \texttt{CLAUDE.md} — Project context/instructions
  \item \texttt{\textasciitilde/.claude/CLAUDE.md} — Global memory
\end{itemize}

\vspace{2mm}
\renewcommand{\arraystretch}{1.15}
\begin{tabularx}{\linewidth}{@{}>{\ttfamily\small}l>{\raggedright\arraybackslash}X@{}}
\normalfont\textbf{Setting} & \textbf{Purpose}\\
\hline
theme & dark / light / system\\
model & Default model to use\\
permissions & Tool approval rules\\
\end{tabularx}
\renewcommand{\arraystretch}{1}

\section*{\textcolor{Cerulean}{7. MCP SERVERS}}
Model Context Protocol extends Claude's capabilities:

\renewcommand{\arraystretch}{1.15}
\begin{tabularx}{\linewidth}{@{}>{\raggedright\arraybackslash}X@{}}
\textbf{Common MCP servers:}\\
\hline
\texttt{@anthropic-ai/mcp-server-memory} — Persistent memory\\
\texttt{@anthropic-ai/mcp-server-fetch} — Enhanced web fetch\\
\texttt{@anthropic-ai/mcp-server-github} — GitHub integration\\
\texttt{@anthropic-ai/mcp-server-postgres} — Database access\\
\end{tabularx}
\renewcommand{\arraystretch}{1}

\vspace{2mm}
\begin{tcolorbox}[colback=Cerulean!8, colframe=Cerulean, boxsep=2pt, arc=1pt, fontupper=\small]
Configure in \texttt{settings.json} under \texttt{"mcpServers"} key.
\end{tcolorbox}

\section*{\textcolor{Violet}{8. HOOKS}}
Automate actions on Claude events:

\renewcommand{\arraystretch}{1.15}
\begin{tabularx}{\linewidth}{@{}>{\ttfamily\small}l>{\raggedright\arraybackslash}X@{}}
\normalfont\textbf{Hook} & \textbf{Triggers When}\\
\hline
PreToolUse & Before tool execution\\
PostToolUse & After tool execution\\
Notification & Claude sends notification\\
Stop & Conversation ends\\
\end{tabularx}
\renewcommand{\arraystretch}{1}

\vspace{2mm}
\begin{tcolorbox}[colback=Violet!8, colframe=Violet, boxsep=2pt, arc=1pt, fontupper=\small\itshape]
Hooks can block, modify, or log tool usage.
\end{tcolorbox}

\section*{\textcolor{OliveGreen}{9. EFFECTIVE PROMPTING}}
\renewcommand{\arraystretch}{1.15}
\begin{tabularx}{\linewidth}{@{}>{\raggedright\arraybackslash}X@{}}
\textbf{Best practices for Claude Code:}\\
\hline
Be specific about file paths and changes\\
Provide context in CLAUDE.md for project conventions\\
Use "don't" to prevent unwanted behaviors\\
Ask Claude to explain before implementing\\
Request tests alongside new features\\
Use \texttt{@file.txt} to include file contents\\
\end{tabularx}
\renewcommand{\arraystretch}{1}

\vspace{2mm}
\begin{tcolorbox}[colback=OliveGreen!8, colframe=OliveGreen, boxsep=2pt, arc=1pt, fontupper=\small]
\textbf{Example CLAUDE.md:}\\
\texttt{Use TypeScript with strict mode.}\\
\texttt{Run tests with: npm test}\\
\texttt{Don't modify package.json without asking.}
\end{tcolorbox}

\section*{\textcolor{RawSienna}{10. COMMON WORKFLOWS}}
\renewcommand{\arraystretch}{1.15}
\begin{tabularx}{\linewidth}{@{}>{\raggedright\arraybackslash}X@{}}
\textbf{Typical tasks and how to request them:}\\
\hline
"Fix the bug in auth.ts line 42"\\
"Add unit tests for the UserService class"\\
"Refactor this function to use async/await"\\
"Explain how the payment flow works"\\
"Review my changes before I commit"\\
"Find all usages of deprecated API"\\
"Create a new React component for..."\\
\end{tabularx}
\renewcommand{\arraystretch}{1}

\vspace{2mm}
\begin{tcolorbox}[colback=RawSienna!8, colframe=RawSienna, boxsep=2pt, arc=1pt, fontupper=\small\itshape]
Claude reads your codebase for context — point it to relevant files.
\end{tcolorbox}

\section*{\textcolor{Gray}{11. TIPS \& TRICKS}}
\renewcommand{\arraystretch}{1.15}
\begin{tabularx}{\linewidth}{@{}>{\raggedright\arraybackslash}X@{}}
\textbf{Power user techniques:}\\
\hline
Use \texttt{claude -c} to continue where you left off\\
Pipe error output: \texttt{npm test 2>\&1 | claude}\\
Run headless in CI with \texttt{-p} flag\\
Check costs with \texttt{/cost} during long sessions\\
Use \texttt{/compact} to reduce context usage\\
Create project-specific CLAUDE.md files\\
Set up hooks to auto-format on file changes\\
\end{tabularx}
\renewcommand{\arraystretch}{1}

\vfill

\end{multicols}
\end{document}
